\documentclass[../../main.tex]{subfiles}

\begin{document}

\section{Kantorovich 对偶}\label{section:ch19-kantorovich}

Riesz 表示定理告诉我们, 对 $1\leqslant p<\infty$ 且 $\sigma$ 有限测度空间, $L^{p}(X,\mu)$ 的对偶空间等距同构于 $L^{q}(X,\mu)$, 其中 $1/p+1/q=1$. 对于 $p=\infty$, $L^{\infty}(X,\mu)$ 的对偶空间通常比 $L^{1}(X,\mu)$ 大. Kantorovich 表示定理告诉我们, $L^{\infty}(X,\mu)$ 的对偶空间可以用关于 $\mu$ 绝对连续的有限带号测度来刻画.

\begin{theorem}[Kantorovich 表示定理]\label{theorem:ch19-kantorovich}
令 $(X,\mathcal{M},\mu)$ 为 $\sigma$ 有限测度空间. $L^{\infty}(X,\mu)$ 上的每个有界线性泛函 $S$ 都具有以下形式: 存在 $\mathcal{M}$ 上的有限带号测度 $\nu$, 它关于 $\mu$ 绝对连续, 使得
\begin{align}
\text{对所有 }f\in L^{\infty}(X,\mu),\quad S(f)=\int_{X}f\,d\nu.\label{align:ch19-kantorovich}
\end{align}
此外, $\|S\|=|\nu|(X)$.
\end{theorem}

\begin{proof}
定义 $\mathcal{M}$ 上的集函数 $\nu$ 为对 $E\in\mathcal{M}$, $\nu(E)=S(\chi_{E})$. 则 $\nu$ 是有限可加的. 为证明 $\nu$ 是可数可加的, 令 $\{E_{k}\}$ 为不交可测集族. 则 $\chi_{\bigcup_{k=1}^{n}E_{k}}$ 在 $L^{\infty}$ 中收敛于 $\chi_{\bigcup E_{k}}$, 因此
\begin{align*}
\nu\left(\bigcup_{k=1}^{\infty}E_{k}\right)=\lim_{n\to\infty}\nu\left(\bigcup_{k=1}^{n}E_{k}\right)=\sum_{k=1}^{\infty}\nu(E_{k}).
\end{align*}
若 $\mu(E)=0$, 则 $\chi_{E}=0$ a.e., 因此 $\nu(E)=0$. 于是 $\nu\ll\mu$. 由 Radon-Nikodym 定理, 存在 $g\in L^{1}(X,\mu)$, 使得 $\nu(E)=\int_{E}g\,d\mu$. 然而, 一般来说 $g$ 不一定在 $L^{\infty}$ 中; 泛函 $S$ 作用在 $L^{\infty}$ 上, 因此 $\nu$ 的全变差测度 $|\nu|$ 是有限的. 对简单函数 $\psi$, $S(\psi)=\int\psi\,d\nu$, 通过延拓到 $L^{\infty}$ 得到\eqref{align:ch19-kantorovich}. 范数等式从 $|\nu|(X)=\sup\{|S(f)|:\|f\|_{\infty}\leqslant 1\}$ 得到.
\end{proof}

我们称 $L^{\infty}(X,\mu)$ 的对偶空间为 $ba(\mathcal{M},\mu)$, 即关于 $\mu$ 绝对连续的有限带号测度空间. 注意这与 $L^{1}(X,\mu)$ 不同: $L^{1}(X,\mu)$ 是 $L^{\infty}(X,\mu)$ 的对偶空间的子空间.

\begin{exercise}\label{exercise:ch19-09}
验证: 若 $\mu$ 是计数测度在自然数集上, 则 $L^{\infty}(\mathbf{N},\mu)=l^{\infty}$, 它的对偶空间是 $ba$, 即有界有限可加测度空间, 它比 $l^{1}$ 大.
\end{exercise}

\end{document}